\section{Appendix B. Implementation Details}

\subsection{B.1 Core Module}

\begin{lstlisting}
class ChelaSlayAttention(BaseAttention):
    def \_\_init\_\_(self, config, rank=64, feature\_dim=128, gate\_init\_bias=5.0, ...):
        super().\_\_init\_\_(config, rank)
        
        self.dim = ...
        self.num\_heads = ...
        self.head\_dim = ...
        self.feature\_dim = feature\_dim
        
        # LoRA projections
        self.W\_q = LoRALinear(...)
        self.W\_k = LoRALinear(...)
        self.W\_v = LoRALinear(...)
        self.out\_proj = LoRALinear(...)
        
        # Yat-Kernel RFF
        self.register\_buffer("random\_proj", 
                             torch.randn(self.head\_dim, feature\_dim) * 0.02)
        
        # Adaptive gate
        self.gate = nn.Linear(self.dim, self.num\_heads)
        nn.init.zeros\_(self.gate.weight)
        self.gate.bias.data.fill\_(gate\_init\_bias)
\end{lstlisting}

\subsection{B.2 Spherical Normalization}

\begin{lstlisting}
def \_spherical\_normalize(self, x):
    return F.normalize(x, p=2, dim=-1, eps=1e-6)
\end{lstlisting}

\subsection{B.3 Yat-Kernel Feature Map}

\begin{lstlisting}
def \_yat\_feature\_map(self, x):
    if self.random\_proj.shape[0] != x.shape[-1]:
        self.random\_proj = torch.randn(
            x.shape[-1], self.feature\_dim, 
            device=x.device, dtype=x.dtype
        ) * 0.02
    proj = torch.matmul(x, self.random\_proj.to(device=x.device, dtype=x.dtype))
    return torch.cat([torch.cos(proj), torch.sin(proj)], dim=-1) * (
        1.0 / (self.feature\_dim ** 0.25)
    )
\end{lstlisting}

\subsection{B.4 Causal Linear Attention}

\begin{lstlisting}
# Outer product
outer\_prod = k\_feat.float().unsqueeze(-1) * v\_masked.float().unsqueeze(-2)
# Cumsum
S = torch.cumsum(outer\_prod, dim=2)
Z = torch.cumsum(k\_feat.float(), dim=2)

num = torch.matmul(q\_feat.float().unsqueeze(-2), S).squeeze(-2)
denom = (q\_feat.float() * Z).sum(dim=-1, keepdim=True) + 1e-5
linear\_out = (num / denom.clamp(min=1e-5)).to(x.dtype)
\end{lstlisting}

\subsection{B.5 Bidirectional Linear Attention}

\begin{lstlisting}
KV = torch.matmul(k\_feat.transpose(-2, -1).float(), v\_masked.float())
Z = k\_feat.float().sum(dim=2, keepdim=True)
num = torch.matmul(q\_feat.float(), KV)
denom = (q\_feat.float() * Z).sum(dim=-1, keepdim=True) + 1e-5
linear\_out = (num / denom.clamp(min=1e-5)).to(x.dtype)
\end{lstlisting}

\subsection{B.6 Gate Fusion}

\begin{lstlisting}
g = torch.sigmoid(self.gate(x)).transpose(1, 2).unsqueeze(-1)
fused = g * dense\_out + (1.0 - g) * linear\_out
\end{lstlisting}

\subsection{B.7 Common Pitfalls}

\subsubsection{B.7.1 Memory OOM in Causal Mode}

The cumsum intermediate is large. Reduce batch size or feature\_dim.

\subsubsection{B.7.2 Random Proj Not Moved to Device}

Ensure \texttt{.to(device=x.device)} when the buffer is on a different device.

\subsubsection{B.7.3 Missing NaN Protection}

Always \texttt{torch.nan\_to\_num} after softmax.

\subsubsection{B.7.4 Gate Init Too Small}

$b_g < 3$ may allow the linear stream to dominate early, causing instability.

\subsubsection{B.7.5 Wrong Nested Encoder Handling}

For RoBERTa, \texttt{is\_nested\_encoder=True} skips \texttt{out\_proj} because the outer 
\texttt{attention.output.dense} handles it.

\subsection{B.8 Feature Dim Auto-Adjustment}

The \texttt{random\_proj} buffer is registered with \texttt{head\_dim} rows. If the 
pretrained attention has a different head\_dim than expected, the buffer is 
recreated on the fly. This is safe but inefficient.

